%%%%%%%%%%%%%%%%%%%%%%%%%%%%%%%%%%%%%%%%%%%%%%%%%%%%%%%%%%%%%%%%%%%%%%
% Lớp 9 | Chương V. Đường tròn | Độ dài của cung tròn. Diện tích hình quạt tròn và hình vành khuyên
% Dạng 9H5?15-1: Độ dài đường tròn, cung tròn
%%%%%%%%%%%%%%%%%%%%%%%%%%%%%%%%%%%%%%%%%%%%%%%%%%%%%%%%%%%%%%%%%%%%%%


% nguồn: dan_tu_AI (nhập 14/11/2026)
\begin{ex}%[9H5N15-1]%[ma:2514ce01]
	Trên đường tròn bán kính $R$, độ dài $l$ của cung tròn có số đo $n^\circ$ được tính bởi công thức
	\choice
	{$l=\dfrac{\pi Rn}{360}$}
	{$l=\dfrac{\pi R^2n}{180}$}
	{\True $l=\dfrac{\pi Rn}{180}$}
	{$l=\dfrac{\pi R^2n}{360}$}
	\loigiai{
		Độ dài cung tròn $n^\circ$ của đường tròn bán kính $R$ là $l=\dfrac{\pi Rn}{180}$.
	}
\end{ex}

% nguồn: dan_tu_AI (nhập 14/11/2026)
\begin{ex}%[9H5H15-1]%[ma:5db94c54]
	Một bánh xe có đường kính $70\text{ cm}$. Khi bánh xe lăn đủ $10$ vòng thì xe đi được quãng đường bao nhiêu mét (lấy $\pi\approx 3{,}14$)?
	\choice
	{\True $21{,}98\text{ m}$}
	{$43{,}96\text{ m}$}
	{$10{,}99\text{ m}$}
	{$2{,}198\text{ m}$}
	\loigiai{
		Chu vi bánh xe: $C=\pi d\approx 3{,}14\cdot 70=219{,}8\text{ (cm)}$.\\
		Quãng đường sau $10$ vòng: $10\cdot 219{,}8=2198\text{ (cm)}=21{,}98\text{ (m)}$.
	}
\end{ex}

% nguồn: dan_tu_AI (nhập 14/11/2026)
\begin{ex}%[9H5H15-1]%[ma:9f678e5d]
	Cho đường tròn $(O;6\text{ cm})$ và hai điểm $A$, $B$ thuộc đường tròn sao cho số đo cung nhỏ $AB$ bằng $120^\circ$. Các khẳng định sau đúng hay sai?
	\choiceTF
	{\True Chu vi đường tròn $(O)$ bằng $12\pi\text{ cm}$}
	{\True Độ dài cung nhỏ $AB$ bằng $4\pi\text{ cm}$}
	{Độ dài cung lớn $AB$ bằng $6\pi\text{ cm}$}
	{\True Một cung của đường tròn $(O)$ có độ dài $3\pi\text{ cm}$ thì có số đo bằng $90^\circ$}
	\loigiai{
		\begin{itemchoice}
			\itemch Đúng. $C=2\pi R=2\pi\cdot 6=12\pi\text{ (cm)}$.
			\itemch Đúng. $l=\dfrac{\pi\cdot 6\cdot 120}{180}=4\pi\text{ (cm)}$.
			\itemch Sai. Cung lớn $AB$ có số đo $240^\circ$, độ dài $\dfrac{\pi\cdot 6\cdot 240}{180}=8\pi\text{ (cm)}$.
			\itemch Đúng. Từ $l=\dfrac{\pi Rn}{180}$ suy ra $n=\dfrac{180l}{\pi R}=\dfrac{180\cdot 3\pi}{6\pi}=90$.
		\end{itemchoice}
	}
\end{ex}

% nguồn: dan_tu_AI (nhập 18/11/2026)
\begin{bt}%[9H5V15-1]%[ma:d7f113cf]
	Kim phút của một chiếc đồng hồ treo tường dài $12\text{ cm}$. Hỏi trong thời gian $25$ phút, đầu mút của kim phút vạch nên một cung tròn có độ dài bằng bao nhiêu? (Lấy $\pi \approx 3{,}14$ và làm tròn kết quả đến hàng phần mười).
	\loigiai{
		Kim phút của đồng hồ quét một vòng tròn trọn vẹn (tương ứng với góc ở tâm $360^\circ$) trong thời gian $60$ phút.\\
		Góc ở tâm mà kim phút vạch được trong vòng $1$ phút là: $\dfrac{360^\circ}{60} = 6^\circ$.\\
		Trong thời gian $25$ phút, kim phút quay được một góc ở tâm có số đo là: $n = 25 \cdot 6^\circ = 150^\circ$.\\
		Đầu mút của kim phút chuyển động tạo thành một cung tròn có bán kính $R = 12\text{ cm}$ (đúng bằng chiều dài kim phút).\\
		Độ dài cung tròn mà đầu mút kim phút vạch được là:\\
		$l = \dfrac{\pi R n}{180} = \dfrac{\pi \cdot 12 \cdot 150}{180} = \dfrac{1800\pi}{180} = 10\pi\text{ (cm)}$.\\
		Sử dụng $\pi \approx 3{,}14$, ta tính được: $l \approx 10 \cdot 3{,}14 = 31{,}4\text{ (cm)}$.\\
		Vậy trong $25$ phút, đầu mút của kim phút vạch nên một cung tròn dài khoảng $31{,}4\text{ cm}$.
	}
\end{bt}

% nguồn: 2026-11-21_lop9_chuong5_bai14_17 (nhập 21/11/2026)
\begin{ex}%[9H5N15-1]%[ma:5d4d672a]
% nguon: SGK Toán 9 KNTT tập 1 – Bài 15, Độ dài của cung tròn. Diện tích hình quạt tròn và hình vành khuyên
	Độ dài cung tròn $60^\circ$ của đường tròn bán kính $R = 9\text{ cm}$ bằng
	\choice
	{$\dfrac{3\pi}{2}\text{ cm}$}
	{\True $3\pi\text{ cm}$}
	{$6\pi\text{ cm}$}
	{$9\pi\text{ cm}$}
	\loigiai{
		Độ dài cung tròn $60^\circ$ bán kính $R = 9\text{ cm}$ là $l = \dfrac{n}{180}\pi R = \dfrac{60}{180}\pi \cdot 9 = 3\pi\text{ cm}$.
	}
\end{ex}

% nguồn: 2026-11-21_lop9_chuong5_bai14_17 (nhập 21/11/2026)
\begin{ex}%[9H5H15-1]%[ma:cbb00f9d]
% nguon: SGK Toán 9 KNTT tập 1 – Bài 15, Độ dài của cung tròn. Diện tích hình quạt tròn và hình vành khuyên
	Cho đường tròn $(O; 12\text{ cm})$ và cung nhỏ $\overset{\frown}{AB}$ có số đo $90^\circ$. Xét tính đúng, sai của các khẳng định sau:
	\choiceTF
	{\True Chu vi của đường tròn $(O)$ bằng $24\pi\text{ cm}$}
	{\True Độ dài cung nhỏ $\overset{\frown}{AB}$ bằng $6\pi\text{ cm}$}
	{Độ dài cung lớn $\overset{\frown}{AB}$ bằng $16\pi\text{ cm}$}
	{\True Chu vi của hình quạt tròn $OAB$ (ứng với cung nhỏ $\overset{\frown}{AB}$), lấy $\pi \approx 3{,}14$ và làm tròn đến hàng phần mười, bằng $42{,}8\text{ cm}$}
	\loigiai{
		\begin{itemchoice}
			\itemch Đúng. Chu vi đường tròn là $C = 2\pi R = 2\pi \cdot 12 = 24\pi\text{ cm}$.
			\itemch Đúng. Độ dài cung nhỏ $\overset{\frown}{AB}$ là $l = \dfrac{90}{180}\pi \cdot 12 = 6\pi\text{ cm}$.
			\itemch Sai. Độ dài cung lớn $\overset{\frown}{AB}$ là $24\pi - 6\pi = 18\pi\text{ cm}$.
			\itemch Đúng. Chu vi hình quạt tròn $OAB$ bằng $2R + l = 2 \cdot 12 + 6\pi \approx 24 + 6 \cdot 3{,}14 = 42{,}84 \approx 42{,}8\text{ cm}$.
		\end{itemchoice}
	}
\end{ex}

% nguồn: 2026-11-21_lop9_chuong5_bai14_17 (nhập 21/11/2026)
\begin{bt}%[9H5V15-1]%[ma:fc336a9a]
% nguon: SGK Toán 9 KNTT tập 1 – Bài 15, Độ dài của cung tròn. Diện tích hình quạt tròn và hình vành khuyên
	Một bánh xe đạp có bán kính $35\text{ cm}$. Khi bánh xe lăn không trượt được $10$ vòng trên một đoạn đường thẳng thì xe đi được quãng đường dài bao nhiêu mét? (Lấy $\pi \approx 3{,}14$, làm tròn kết quả đến hàng đơn vị).
	\shortans{22}
	\loigiai{
		Đổi $R = 35\text{ cm} = 0{,}35\text{ m}$.\\
		Chu vi của bánh xe là $C = 2\pi R \approx 2 \cdot 3{,}14 \cdot 0{,}35 = 2{,}198\text{ m}$.\\
		Quãng đường xe đi được sau $10$ vòng lăn là $10 \cdot 2{,}198 = 21{,}98 \approx 22\text{ m}$.
	}
\end{bt}

% nguồn: 2026-11-21_lop9_chuong5_bai14_15_bs (nhập 21/11/2026)
\begin{ex}%[9H5N15-1]%[ma:4bbb8a5b]
% nguon: SGK Toán 9 KNTT tập 1 – Bài 15, Độ dài của cung tròn. Diện tích hình quạt tròn và hình vành khuyên
	Công thức tính độ dài $l$ của cung tròn $n^\circ$ trên đường tròn bán kính $R$ là
	\choice
	{\True $l = \dfrac{n}{180}\pi R$}
	{$l = \dfrac{n}{360}\pi R$}
	{$l = \dfrac{n}{180}\pi R^2$}
	{$l = \dfrac{n}{360}\pi R^2$}
	\loigiai{
		Trên đường tròn bán kính $R$, độ dài cung tròn có số đo $n^\circ$ được tính theo công thức $l = \dfrac{n}{180}\pi R$.
	}
\end{ex}

% nguồn: 2026-11-21_lop9_chuong5_bai14_15_bs (nhập 21/11/2026)
\begin{ex}%[9H5N15-1]%[ma:c9fa2903]
% nguon: SGK Toán 9 KNTT tập 1 – Bài 15, Độ dài của cung tròn. Diện tích hình quạt tròn và hình vành khuyên
	Chu vi của đường tròn có đường kính $d = 10\text{ cm}$ bằng
	\choice
	{$5\pi\text{ cm}$}
	{\True $10\pi\text{ cm}$}
	{$20\pi\text{ cm}$}
	{$25\pi\text{ cm}$}
	\loigiai{
		Chu vi đường tròn có đường kính $d = 10\text{ cm}$ là $C = \pi d = 10\pi\text{ cm}$.
	}
\end{ex}

% nguồn: 2026-11-21_lop9_chuong5_bai14_15_bs (nhập 21/11/2026)
\begin{ex}%[9H5H15-1]%[ma:f680d67e]
% nguon: SGK Toán 9 KNTT tập 1 – Bài 15, Độ dài của cung tròn. Diện tích hình quạt tròn và hình vành khuyên
	Cho đường tròn $(O; 18\text{ cm})$ và cung nhỏ $\overset{\frown}{MN}$ có số đo $40^\circ$. Xét tính đúng, sai của các khẳng định sau:
	\choiceTF
	{\True Độ dài đường tròn $(O)$ bằng $36\pi\text{ cm}$}
	{\True Độ dài cung nhỏ $\overset{\frown}{MN}$ bằng $4\pi\text{ cm}$}
	{Độ dài cung lớn $\overset{\frown}{MN}$ bằng $30\pi\text{ cm}$}
	{\True Nếu giữ nguyên bán kính và tăng số đo cung lên $120^\circ$ thì độ dài cung mới gấp $3$ lần độ dài cung nhỏ $\overset{\frown}{MN}$}
	\loigiai{
		\begin{itemchoice}
			\itemch Đúng. Độ dài đường tròn là $C = 2\pi \cdot 18 = 36\pi\text{ cm}$.
			\itemch Đúng. Độ dài cung nhỏ $\overset{\frown}{MN}$ là $l = \dfrac{40}{180}\pi \cdot 18 = 4\pi\text{ cm}$.
			\itemch Sai. Độ dài cung lớn $\overset{\frown}{MN}$ là $36\pi - 4\pi = 32\pi\text{ cm}$.
			\itemch Đúng. Với bán kính không đổi, độ dài cung tròn tỉ lệ thuận với số đo cung; vì $120^\circ = 3 \cdot 40^\circ$ nên độ dài cung mới gấp $3$ lần.
		\end{itemchoice}
	}
\end{ex}

% nguồn: 2026-11-21_lop9_chuong5_bai14_15_bs (nhập 21/11/2026)
\begin{ex}%[9H5H15-1]%[ma:b2af6259]
% nguon: SGK Toán 9 KNTT tập 1 – Bài 15, Độ dài của cung tròn. Diện tích hình quạt tròn và hình vành khuyên
	Một cung tròn trên đường tròn bán kính $R = 15\text{ cm}$ có độ dài bằng $5\pi\text{ cm}$. Xét tính đúng, sai của các khẳng định sau:
	\choiceTF
	{\True Chu vi đường tròn chứa cung đó bằng $30\pi\text{ cm}$}
	{\True Số đo của cung tròn đó bằng $60^\circ$}
	{Độ dài dây căng cung tròn đó bằng $10\text{ cm}$}
	{\True Chu vi hình quạt tròn tạo bởi cung đó (lấy $\pi \approx 3{,}14$) bằng $45{,}7\text{ cm}$}
	\loigiai{
		\begin{itemchoice}
			\itemch Đúng. Chu vi đường tròn là $C = 2\pi \cdot 15 = 30\pi\text{ cm}$.
			\itemch Đúng. Từ $l = \dfrac{n}{180}\pi R \Rightarrow 5\pi = \dfrac{n}{180}\pi \cdot 15 \Rightarrow n = 60^\circ$.
			\itemch Sai. Góc ở tâm bằng $60^\circ$ nên tam giác tạo bởi tâm và hai mút cung là tam giác đều cạnh $15\text{ cm}$, suy ra dây dài $15\text{ cm}$.
			\itemch Đúng. Chu vi hình quạt tròn bằng $2R + l = 30 + 5\pi \approx 30 + 5 \cdot 3{,}14 = 45{,}7\text{ cm}$.
		\end{itemchoice}
	}
\end{ex}

% nguồn: 2026-11-21_lop9_chuong5_bai14_15_bs (nhập 21/11/2026)
\begin{bt}%[9H5V15-1]%[ma:4bd672f7]
% nguon: SGK Toán 9 KNTT tập 1 – Bài 15, Độ dài của cung tròn. Diện tích hình quạt tròn và hình vành khuyên
	Một chiếc đu quay có bán kính $20\text{ m}$. Khi đu quay quay một góc $90^\circ$ thì một cabin trên vành đu quay di chuyển được quãng đường dài bao nhiêu mét? (Lấy $\pi \approx 3{,}14$, làm tròn kết quả đến hàng phần mười).
	\shortans{31,4}
	\loigiai{
		Quãng đường cabin di chuyển là độ dài cung tròn $90^\circ$ bán kính $R = 20\text{ m}$:
		$$l = \dfrac{90}{180}\pi \cdot 20 = 10\pi \approx 10 \cdot 3{,}14 = 31{,}4\text{ m}.$$
	}
\end{bt}

% nguồn: 2026-11-21_lop9_chuong5_bai14_15_bs (nhập 21/11/2026)
\begin{bt}%[9H5V15-1]%[ma:791b7c63]
% nguon: SGK Toán 9 KNTT tập 1 – Bài 15, Độ dài của cung tròn. Diện tích hình quạt tròn và hình vành khuyên
	Một cung tròn có số đo $72^\circ$ và độ dài bằng $6\pi\text{ cm}$. Tính bán kính của đường tròn chứa cung đó (đơn vị cm).
	\shortans{15}
	\loigiai{
		Áp dụng công thức $l = \dfrac{n}{180}\pi R$, ta có:
		$$6\pi = \dfrac{72}{180}\pi R \Rightarrow \dfrac{2}{5}R = 6 \Rightarrow R = 15\text{ cm}.$$
	}
\end{bt}
